\chapter{La curiosidad empieza en tierra}
\chapterquote{Aprender comienza cuando una pregunta reemplaza una certeza.}{Nota editorial}
\opening{Una cabina pequeña puede abrir preguntas enormes.}

Este es un capítulo de demostración, no un procedimiento de vuelo. Un libro de divulgación puede alternar relatos personales, observaciones y explicaciones técnicas sin convertirse en un manual operacional.

Aquí aparece una nota al margen\sidenote{Las notas marginales son ideales para una definición corta, un dato o una referencia que no interrumpa el relato.}. La composición está inspirada en la tradición editorial de Edward Tufte, pero construida sobre \texttt{memoir}.

\section{La diferencia entre mirar y comprender}

La pregunta interesante no es solo qué hace una pieza, sino qué problema resuelve. Esta clase de libro se beneficia de una lectura continua y de ilustraciones puntuales.

\begin{insightbox}[Idea para recordar]
Una buena nota visual debe aclarar el contenido principal, no competir con él.
\end{insightbox}

\margintext{Una idea lateral puede leerse por separado, como segunda capa de contenido.}Un texto más largo permite comprobar el uso de notas y el espaciado entre líneas. La plantilla acepta capítulos independientes y bibliografía administrada por BibLaTeX. Una referencia de ejemplo puede citarse con \textcite{tufte2001}.

\section{Una imagen que explique}

\aviamarginfigure{assets/wing.pdf}{Esquema simplificado, dibujado para demostrar una figura marginal.}La figura lateral no necesita ocupar una página completa. Puede acompañar una explicación breve sin interrumpir el párrafo.

\begin{fullwidth}
\begin{insightbox}[Lectura extendida]
Este recuadro ocupa el ancho del texto y el margen lateral. Úsalo para conceptos relevantes, ilustraciones, resúmenes o fotografías panorámicas. En libros de impresión a doble cara es importante probar el diseño de márgenes interiores y exteriores.
\end{insightbox}
\end{fullwidth}
